\section{Related Work}
\label{sec:related}

\textbf{Paired supervision and data acquisition.}
Early reflection-removal benchmarks and perceptual training methods emphasized
the difficulty of separating overlapping layers from a single
image~\cite{wan2017benchmarking,zhang2018perceptual}. More recent real-image
collections, including RRW and OpenRR, improve the relevance of training
data to photographs captured outside synthetic compositing
pipelines~\cite{zhu2024rrw,openrr2025,cai2026ntire}.
Their construction also highlights the importance of how a clean reference
is acquired. Removing a reflecting surface can change refraction and alignment,
whereas covering it can change illumination and camera response.
RAW-domain acquisition offers another route to reducing image-processing
mismatch~\cite{kee2025raw}. Our study focuses on a related but distinct failure of pairing:
partial camera-side occlusion on tinted automotive glass.
We retain the optical tint in the target and examine where changes introduced
by the acquisition procedure affect the loss.

\textbf{Reflection models and restoration networks.}
Classical methods constrain the two layers through image
statistics~\cite{wan2017benchmarking,wan2022crrn}.
Learning-based work extends the model with cascaded
refinement~\cite{li2020ibcln}, non-additive interactions~\cite{hu2023dsrnet},
absorption-aware imaging~\cite{zheng2021absorption}, and reversible
feature separation~\cite{zhao2025rdnet,hu2024dsit}.
Recent diffusion-based approaches provide stronger image
priors~\cite{hong2024ldiffer,promptrr2024,li2026gensirr},
while lightweight architectures address inference
cost~\cite{lightrr2026}.
These methods seek improved restoration under their respective training
distributions. Our question is complementary: when a paired target
has local and global photometric ambiguities, which aspects of the
supervision remain useful under a fixed restoration architecture?

\textbf{Spatial priors and photometric conditions.}
Reflection-location predictors guide restoration by estimating where
reflected content is present. RRW uses a gradient-comparison pseudo-label
and a learned guidance map~\cite{zhu2024rrw}; location-aware and
reflection-guided methods explore related spatial
signals~\cite{lanet2021,ragnet2023}.
A camera-side occlusion map is different: it describes a change in
the reflection-generating process, not the support of the visible
reflection. We compare these signals and separately test whether
more accurate reflection-location labels translate into a better
reconstruction objective.

Photometric conditioning has close parallels in image
enhancement, where learned curves, reflectance--illumination
decompositions, and restoration prompts provide ways to control
appearance~\cite{guo2020zerodce,ma2022sci,wu2022uretuinex,cai2023retinexformer}.
Such methods motivate an operational question: does changing an
explicit condition actually change the output as intended?
We use an intervention on the conditioning value to measure this
response directly, without treating a predicted brightness scalar
as an independent observation.
